% Part: sets-functions-relations
% Chapter: arithmetization
% Section: cauchy

\documentclass[../../../include/open-logic-section]{subfiles}

\begin{document}

\olfileid{sfr}{arith}{cauchy}

\olsection{Lampiran: Wilangan Real minangka Rerangken Cauchy}

Ing \olref[cuts]{sec}, kita mbangun wilangan real minangka potongan
Dedekind. Ing perangan iki, kita njlentrehake pambangunan alternatif.
Pambangunan iki adhedhasar definisi Cauchy tumrap apa kang saiki diarani
rerangken Cauchy; nanging panganggone definisi iki kanggo \emph{mbangun}
wilangan real asale saka panulis abad kaping 19 liyane, mligine Weierstrass,
Heine, M\'{e}ray, lan Cantor. (Kanggo sejarah kang becik, delengen
\citeauthor{OConnorRobertson:RN} \citeyear{OConnorRobertson:RN}.)

Sadurunge tekan abad kaping 19, prayoga yen kita nggatekake Simon Stevin
(1548--1620). Ringkese, Stevin mangerteni yen saben wilangan real bisa
dipikirake liwat ekspansi desimale. Mula malah wilangan irasional kaya
$\sqrt{2}$ duwe ekspansi desimal kang becik, kang diwiwiti:
\[
	1.41421356237\ldots
\]
Gampang banget nggambarake ekspansi desimal ing teori himpunan: cukup
anggepen minangka fungsi $d \colon \Nat \to \Nat$, kanthi $d(n)$ dadi
panggonan desimal kaping $n$ kang digatekake. Banjur perlu owah-owahan
sethithik kanggo nangani bagean wilangan real sadurunge tandha desimal
(ing kene mung $1$). Uga perlu owah-owahan liyane, yaiku relasi
ekuivalensi, kanggo njamin yen, umpamane, $0.999\ldots = 1$. Nanging
ora angel menehi pambangunan wilangan real kang ketat babar pisan
kanthi cara Stevin ing teori himpunan.

Stevin dudu punjering rembugan kita. (Kanggo katrangan luwih akeh bab
Stevin, delengen \citealt{KatzKatz2012}.) Nanging ana gagasan kang raket.
Tinimbang langsung nganggep ekspansi desimal $\sqrt{2}$, kita bisa
nggatekake \emph{rerangken} panyedhakan rasional marang $\sqrt{2}$ kang
saya tliti, nganggo ekspansi kang saya cetha:
\[
1, 1.4, 1.414, 1.4142, 1.41421,\ldots
\]
Gagasan yen wilangan real bisa dipikirake liwat ``panyedhakan kang saya
becik'' dadi dhasar kanggo rerangken wawasan liyane (padha karo pangamatan
kang digunakake nalika mbangun $\Rat$ saka $\Int$, utawa $\Int$ saka
$\Nat$). Wawasan dhasare yaiku:
\begin{enumerate}
	\item Saben wilangan real bisa ditulis minangka ekspansi desimal
	(kang bisa tanpa wates).
	\item Informasi kang dikodeake ekspansi desimal (kang bisa tanpa wates)
	uga bisa dikodeake kanthi rerangken wilangan rasional.
	\item Rerangken wilangan rasional bisa dipikirake minangka fungsi saka
	$\Nat$ menyang $\Rat$; cukup netepna $f(n)$ minangka wilangan rasional
	kaping $n$ ing rerangken kasebut.
\end{enumerate}
Mesthi wae, ora \emph{sembarang} fungsi saka $\Nat$ menyang $\Rat$ menehi
wilangan real. Umpamane, gatekna fungsi iki:
\[
	f(n) =\begin{cases}
			1 & \text{yen }n\text{ ganjil}\\
			0 &\text{yen }n\text{ genep}
		\end{cases}
\]
Inti masalahe yaiku rerangken $0,1,0,1,0,1,0,\ldots$ katon ora
``nyedhaki'' wilangan real apa wae. Dadi, supaya rerangken kang digatekake
pancen nyedhaki sawenehing wilangan real, kita kudu matesi kawigaten mung
ing rerangken kang duwe \emph{limit}.

Kita wis nemoni gagasan limit ing \olref[his][set][limits]{sec}. Nanging
kita ora bisa nggunakake definisi kang \emph{persis padha} karo definisi
ing kono. Ekspresi ``$(\forall \epsilon>0)$'' ing kono kanthi meneng-meneng
nggunakake kuantifikasi ing wilangan real; lan kita nggatekake limit fungsi
ing wilangan real. Dadi nggunakake definisi kasebut ateges wis nganggep
wilangan real cumawis, kamangka wilangan iku persis kang arep
\emph{dibangun}. Untunge, kita bisa nggunakake gagasan limit kang raket.
\begin{defn}\ollabel{def:CauchySequence}
  Fungsi $f: \Nat \to \Rat$ diarani \emph{rerangken Cauchy} yen lan mung
  yen kanggo saben $\epsilon \in \Rat$ kang positif, kita duwe $(\exists
  \ell \in \Nat)(\forall m, n > \ell)|f(m) - f(n)| < \epsilon$.
\end{defn}

Gagasan umum limit isih padha: yen kita kepengin tingkat katliten tartamtu
(kang diukur nganggo~$\epsilon$), ana ``wilayah'' kang kudu dideleng, yaiku
sembarang masukan luwih gedhe tinimbang~$\ell$. Gampang dingerteni yen
rerangken $1$, $1.4$, $1.414$, $1.4142$, $1.41421$\ldots{} duwe limit: yen
kita kepengin nyedhaki $\sqrt{2}$ kanthi kasalahan kurang saka
$\nicefrac{1}{10^{n}}$, cukup delengen sembarang entri sawise entri kaping
$n$.

Gagasan kang katon cetha yaiku nyebut wilangan real \emph{iku} sembarang
rerangken Cauchy. Nanging kaya ing pambangunan $\Int$ lan $\Rat$, iki
naif banget: kanggo saben wilangan real, ana akeh rerangken Cauchy beda kang
nuduhake wilangan real kasebut. Cara prasaja kanggo ndeleng iki kaya
mangkene. Yen diwenehi rerangken Cauchy~$f$, netepna $g$ minangka fungsi
kang persis padha karo~$f$, kejaba $g(0)\neq f(0)$. Amarga rerangken loro
iku padha ing saben panggonan sawise wilangan kapisan, pungkasane kita
kepengin kandha yen rerangken mau duwe limit kang padha miturut teges ing
\olref{def:CauchySequence}, mula kudu dianggep ``netepake'' wilangan real
kang padha. Dadi, rerangken Cauchy loro iki kudu dianggep wilangan real
kang padha. Cathetan koreksi sumber OLP-0048: terjemahan iki mulihake
tembung panyambung kang ilang ing ukara Inggris ``way to see this as follows''.

Mula kita maneh kudu netepake relasi ekuivalensi ing rerangken Cauchy lan
ngenali wilangan real karo kelas-kelas ekuivalensi. Kaping pisan, kita
mbutuhake gagasan fungsi kang nyedhaki $0$ ing limit. Kanggo sembarang
fungsi $h : \Nat \to \Rat$, kandhanana yen \emph{$h$ nyedhaki $0$} yen lan
mung yen kanggo saben $\epsilon \in \Rat$ kang positif kita duwe $(\exists
\ell \in \Nat)(\forall n > \ell)|h(n)| < \epsilon$.\footnote{Bandhingna
iki karo definisi $\lim_{x \mathord{\rightarrow}\infty}f(x) = 0$ ing
\olref[his][set][limits]{sec}.} Sabanjure, yen $f$ lan $g$ iku fungsi
$\Nat \to \Rat$, netepna $(f-g)(n) = f(n) - g(n)$. Saiki netepna:
\[
	f \Realequiv g \text{ yen lan mung yen $(f-g)$ nyedhaki $0$}.
\]
Kita kudu mriksa yen $\Realequiv$ iku relasi ekuivalensi; lan pancen iya.
Banjur, yen dikarepake, kita bisa netepake wilangan real minangka
kelas-kelas ekuivalensi miturut $\Realequiv$ saka kabeh rerangken Cauchy
saka $\Nat \to \Rat$. Cathetan koreksi sumber OLP-0048: ukara Inggris
sadurunge definisi salah nyebut ``equivalence relations''; obyek kang
nggambarake wilangan real yaiku kelas ekuivalensi.

\begin{prob}
Netepna $f(n) = 0$ kanggo saben $n$. Netepna $g(n) = \frac{1}{(n+1)^2}$.
Tuduhna yen kalorone rerangken Cauchy, lan yen limit fungsi kalorone
yaiku $0$, saengga uga $f \sim_\Real g$.
\end{prob}

Sawise iku, kaya lumrahe kita nulis $\equivrep{f}{\Realequiv}$ kanggo
kelas ekuivalensi kang duwe $f$ minangka !!a{element}. Nanging supaya
gampang diwaca, sabanjure kita ngilangi subskrip lan mung nulis
$\equivrep{f}{}$. Kita uga netepake yen kanggo saben $q \in \Rat$, kita
duwe $q_{\Real} = \equivrep{c_{q}}{}$, kanthi $c_{q}$ minangka fungsi
konstan $c_q(n) = q$ kanggo kabeh $n \in \Nat$. Banjur kita netepake
relasi lan operasi dhasar ing wilangan real, umpamane:
\begin{align*}
	\equivrep{f}{} + \equivrep{g}{} &= 	\equivrep{(f + g)}{} \\
	\equivrep{f}{} \times 	\equivrep{g}{} &= \equivrep{(f \times g)}{}
\end{align*}
kanthi $(f + g)(n) = f(n) + g(n)$ lan $(f \times g)(n) = f(n) \times
g(n)$. Mesthi wae, kita uga kudu mriksa yen $(f + g)$, $(f-g)$, lan
$(f\times g)$ iku rerangken Cauchy nalika $f$ lan $g$ rerangken Cauchy;
lan pancen iya, dene buktine dipasrahake marang pamaca.

Pungkasan, kita netepake gagasan tatanan. Kandhanana yen
$\equivrep{f}{}$ iku \emph{positif} yen lan mung yen
$\equivrep{f}{}\neq 0_\Real$ lan $(\exists \ell \in \Nat)(\forall n >
\ell)0 < f(n)$. Banjur kandhanana yen $\equivrep{f}{} < \equivrep{g}{}$
yen lan mung yen $\equivrep{(g - f)}{}$ positif. Kita kudu mriksa yen iki
ditetepake kanthi becik, yaiku ora gumantung marang pilihan fungsi
``wakil'' saka kelas ekuivalensi. Cathetan koreksi sumber OLP-0048:
syarat positif ing sumber nulis nol rasional, nanging kelas kasebut
wilangan real; terjemahan nggunakake nol real kaya ukara sabanjure.
%\begin{align*}
%	\equivrep{f}{} \leq \equivrep{g}{} \text{ yen lan mung yen }&\equivrep{f}{} = \equivrep{g}{} \text{ utawa }\\
%	&(\exists \ell \in \Nat)(\forall n \in \Nat)(\ell < n \lif f(n) \leq g(n))
%\end{align*}
Sawise pamriksa iki, cukup gampang nuduhake yen definisi kasebut menehi
sipat aljabar kang bener, yaiku:
\begin{thm}\ollabel{thm:cauchyorderedfield}
Rerangken-rerangken Cauchy mbentuk lapangan katata. Ing pratelan iki,
kang mbentuk lapangan yaiku kelas-kelas ekuivalensine; rerangken dhewe
digunakake minangka wakil.
\end{thm}

\begin{proof}
Gladhen.
\end{proof}

\begin{prob}
Buktekna yen rerangken-rerangken Cauchy mbentuk lapangan katata, kanthi
teges kelas ekuivalensi kaya kang diterangake ing ndhuwur.
\end{prob}

Luwih angel mbuktekake yen wilangan real kang dibangun kanthi cara iki
nduweni Sipat Kajangkepan, mula kita bakal menehi buktine.

\begin{thm}
Saben himpunan rerangken Cauchy kang ora kosong lan duwe wates ndhuwur
mesthi duwe wates ndhuwur paling cilik. Tatanan lan wates ing pratelan
iki ditrapake liwat kelas ekuivalensi rerangken kasebut.
\end{thm}

\begin{proof}[Gambaran bukti] Upamakna $S$ iku sembarang himpunan rerangken
Cauchy kang ora kosong lan duwe wates ndhuwur. Dadi ana $p \in \Rat$
saengga $p_{\Real}$ dadi wates ndhuwur kanggo $S$. Jupuken $r \in S$;
banjur ana $q \in \Rat$ saengga $q_{\Real} < r$. Dadi yen wates ndhuwur
paling cilik kanggo $S$ ana, wates iku dumunung ing antarane $q_\Real$ lan
$p_\Real$, kalebu loro-lorone. Cathetan sumber OLP-0048: langkah iki
nggunakake kasunyatan Archimedes yen kelas real bisa diwatesi saka ndhuwur
lan ngisor dening salinan rasional; kasunyatan iku durung dibuktekake ing
perangan iki.

Kita bakal saya nyedhaki wates ndhuwur paling cilik kanthi nyedhaki bebarengan
saka ngisor lan ndhuwur. Mligine, kita netepake rong fungsi $f, g \colon
\Nat \to \Rat$, kanthi ancas $f$ nyedhaki wates kasebut saka ndhuwur lan
$g$ nyedhaki wates kasebut\ saka ngisor. Kita miwiti kanthi netepake:
\begin{align*}
	f(0) &= p \\
	g(0) &= q
\end{align*}
Banjur, kanthi $a_n = \frac{f(n) + g(n)}{2}$, netepna:\footnote{Iki
definisi rekursif. Nanging kita \emph{durung} menehi alesan kanggo
nganggep definisi rekursif bisa digunakake.}
\begin{align*}
	f(n+1) &=
	\begin{cases}
		a_n &\text{yen }(\forall h \in S)\equivrep{h}{} \leq (a_n)_\Real\\
		f(n)&\text{yen ora}
	\end{cases}\\
	g(n+1) &=
	\begin{cases}
		a_n &\text{yen }(\exists h \in S)\equivrep{h}{} \geq (a_n)_\Real\\
		g(n) &\text{yen ora}
	\end{cases}
\end{align*}
$f$ lan $g$ kalorone rerangken Cauchy. (Iki bisa dipriksa kanthi cukup
gampang, nanging dipasrahake minangka gladhen.) Gatekna yen fungsi $(f-g)$
nyedhaki $0$, amarga bedane $f$ lan $g$ dadi separo ing saben langkah.
Mula $\equivrep{f}{} = \equivrep{g}{}$.

Kaping pisan kita nuduhake yen $\equivrep{f}{}$ iku wates ndhuwur kanggo
$S$, yaiku\ $(\forall h \in S)\equivrep{h}{} \leq \equivrep{f}{}$.
(Kita bakal nggunakake \olref{thm:cauchyorderedfield} nalika dibutuhake.)
Jupuken $h \in S$ lan upamakna, kanggo bukti kanthi kontradiksi, yen
$\equivrep{f}{} < \equivrep{h}{}$, mula $0_\Real <
\equivrep{(h-f)}{}$. Amarga $f$ iku rerangken Cauchy kang mudhun kanthi
monoton, ana $n \in \Nat$ saengga $\equivrep{(c_{f(n)} - f)}{} <
\equivrep{(h-f)}{}$. Dadi:
\[
	(f(n))_\Real = \equivrep{c_{f(n)}}{} < \equivrep{f}{} + \equivrep{(h-f)}{} = \equivrep{h}{},
\]
kang kosok balen karo kasunyatan yen miturut pambangunan,
$\equivrep{h}{} \leq (f(n))_\Real$.

Sabanjure kita nuduhake yen $\equivrep{f}{} = \equivrep{g}{}$ iku wates
ndhuwur \emph{paling cilik} kanggo $S$. Dadi jupuken sembarang rerangken
Cauchy $j$ lan upamakna $\equivrep{j}{} < \equivrep{g}{}$. Kanthi nalar
kaya ing ndhuwur lan nggunakake kasunyatan yen $g$ \emph{munggah} kanthi
monoton, ana $n \in \Nat$ saengga $\equivrep{j}{} < (g(n))_\Real$.
Nanging miturut pambangunan ana $h \in S$ saengga $(g(n))_\Real \leq
\equivrep{h}{}$, mula $\equivrep{j}{} < \equivrep{h}{}$ lan mulane
$\equivrep{j}{}$ dudu wates ndhuwur kanggo $S$.
\end{proof}

\end{document}
